\documentclass[11pt]{article}
\usepackage[a4paper,margin=1in]{geometry}
\usepackage{amsmath,amssymb}
\usepackage{booktabs,longtable}
\usepackage{graphicx}
\usepackage[hyphens]{url}
\usepackage[colorlinks=true,linkcolor=black,urlcolor=blue,citecolor=blue]{hyperref}
\setlength{\parskip}{0.55em}\setlength{\parindent}{0pt}
\providecommand{\tightlist}{\setlength{\itemsep}{0.2em}\setlength{\parskip}{0pt}}

\title{Joint evaluation of three underselling audits — Paper E1 (2026-09-22)}
\author{\textbf{Amin Abaee}\\ Independent Researcher\\[2pt] ORCID: \href{https://orcid.org/0000-0002-0019-1842}{0000-0002-0019-1842}\\ Email: \href{mailto:amin_abaee@ut.ac.ir}{amin_abaee@ut.ac.ir}}
\date{September 18, 2026}
\begin{document}
\maketitle

\section{Joint evaluation of three audits (deepseek, qwen, grok) on
manuscript underselling --- Paper
E1}\label{joint-evaluation-of-three-audits-deepseek-qwen-grok-on-manuscript-underselling-paper-e1}

\textbf{Date:} 2026-09-22 · \textbf{Scope:}
\texttt{/e1\ underselling.txt} reviews of
\texttt{paperE1\_cod\_forecast\_ladder\_v65} · \textbf{Outcome:}
implemented as v66
(\texttt{paperE1\_cod\_forecast\_ladder\_v66\_generalized.md}, archive +
submission carriers, both 27 pp).

\begin{center}\rule{0.5\linewidth}{0.5pt}\end{center}

\subsection{1. Joint evaluation}\label{joint-evaluation}

The three audits converge on one diagnosis: the manuscript's strongest
contribution is methodological (a pre-registered, benchmark-competitive,
power-aware evaluation protocol), but the framing presents only a
bounded Northern-cod case study. Their remedies converge almost exactly:
modestly generalized title; protocol added to highlights and abstract;
the two named shortcomings elevated as a general critique; a generality
paragraph in Methods; framing sentences in 3.7/3.8; a new Discussion
subsection on broader implications; a generalized closing paragraph in
Conclusions. All three preserve every empirical result and prohibit
overclaims (``persistence is generally superior'', ``structure never
works'').

\textbf{Assessment: the diagnosis is sound and accepted in full.} The
manuscript already contained all the machinery (active-naive-baseline
bridge sentence, frozen-rule documentation, operating-characteristic
simulation, target-robustness re-scoring) --- what was missing is the
explicit claim that this machinery constitutes a reusable protocol. That
is the strongest possible generalization available, because it is fully
backed by material already inside the evaluation; no new result is
claimed.

\subsection{2. Verification of audit claims against v65
(claim-by-claim)}\label{verification-of-audit-claims-against-v65-claim-by-claim}

{\def\LTcaptype{none} % do not increment counter
\begin{longtable}[]{@{}
  >{\raggedright\arraybackslash}p{(\linewidth - 6\tabcolsep) * \real{0.2500}}
  >{\raggedright\arraybackslash}p{(\linewidth - 6\tabcolsep) * \real{0.2500}}
  >{\raggedright\arraybackslash}p{(\linewidth - 6\tabcolsep) * \real{0.2500}}
  >{\raggedright\arraybackslash}p{(\linewidth - 6\tabcolsep) * \real{0.2500}}@{}}
\toprule\noalign{}
\begin{minipage}[b]{\linewidth}\raggedright
\#
\end{minipage} & \begin{minipage}[b]{\linewidth}\raggedright
Audit claim about the manuscript
\end{minipage} & \begin{minipage}[b]{\linewidth}\raggedright
Manuscript evidence (v65 locations)
\end{minipage} & \begin{minipage}[b]{\linewidth}\raggedright
Verdict
\end{minipage} \\
\midrule\noalign{}
\endhead
\bottomrule\noalign{}
\endlastfoot
1 & Abstract limits itself: ``The conclusion applies to this ladder,
this estimator, and this scoring design'' & Abstract, final paragraph,
verbatim & \textbf{Verified} --- but the audits missed that the abstract
already ends with a general sentence (``\ldots cannot be assumed from
structural elaboration alone\ldots{}''). Corrected implementation: add
the \emph{protocol} sentence, not a duplicate lesson sentence. \\
2 & Intro names ``Target Disconnect'' and ``Binary Certification without
Structural Ranking'' as shortcomings & §1, numbered points 1--2,
verbatim emphasis labels & \textbf{Verified} \\
3 & Persistence is treated as a real competitor, while field practice
uses it mainly as a scaling device & §1 bridge paragraph: ``active
competitor rather than a passive scaling denominator''; retention rule
H2 & \textbf{Verified} \\
4 & Scoring rule fixed before scores were read; protocol pre-registered
& Abstract (``scoring rule fixed before the scores were read''); §1
(§3.8 protocol ``pre-registered before any survey score was read'');
§3.7 (``pre-registered simulation experiment'', 200 replicates) &
\textbf{Verified} --- ``pre-registered'' in the title is consistent with
the manuscript's own usage \\
5 & Power analysis: rule highly specific; \textless15\% power for
stock-flow and depensatory dynamics under depleted conditions & Abstract
and §3.7 (97--99\% rejection under persistence null; 97--99\% power
where autonomous signal strong under collapse-window conditions;
\textless15\% for M2/M1b under depletion) & \textbf{Verified} \\
6 & Realized catches supplied to M2--M4; retrospective reconstruction
favours structure ⇒ conservative negative result & §4.5 item 1 (M2, M3,
M4); Abstract ¶2 (``three models receive realised catches\ldots{} makes
the negative result conservative with respect to structure'') &
\textbf{Verified} (grok's ``M2--M4'' exact) \\
7 & §3.8 target robustness: re-scored against raw autumn survey index,
no module retained & §3.8, incl.~+29.8\% vs +17.1\% persistence
advantage; ``What the contrast certifies'' subsection &
\textbf{Verified} \\
8 & Collapse window measures misspecification penalty, not skill (class
cannot generate collapse-then-recovery) & Abstract ¶3; §3.1
collapse-window readings; §3.8 fixed-window readings &
\textbf{Verified} \\
9 & Guardrails: manuscript does not claim best-available model, bears on
skill not status, state-space/age-structured models not excluded &
Abstract (``The result bears on forecast skill, not on stock status; no
tested model is shown to be the best available''); §4.5 item 4 &
\textbf{Verified} --- guardrails preserved in v66 \\
\end{longtable}
}

Every factual premise of the three audits checks out. Nothing in the
audits required the manuscript to gain a new empirical claim.

\subsection{3. Where the audits were weak --- strengthened
jointly}\label{where-the-audits-were-weak-strengthened-jointly}

\begin{enumerate}
\def\labelenumi{\arabic{enumi}.}
\tightlist
\item
  \textbf{duplication risk:} the audits' proposed abstract endings
  repeat the skill-cannot-be-assumed lesson that v65 already carries.
  Resolution: insert one protocol sentence and extend the existing final
  sentence with ``on decision-relevant quantities''.
\item
  \textbf{3.7 redundancy:} §3.7 already opens with ``an empirical
  negative result is only as informative as the statistical power of the
  diagnostic test''. Grok's proposed paragraph would double it.
  Resolution: one closing bookend sentence, framed on the
  instrument/operating-characteristics idea shown in the §3.7 comparison
  table (MASE \textless1: power 0.651/specificity 0.675; IC rule:
  0.509/0.992) --- a column the audits themselves did not cite as
  evidence in their favour.
\item
  \textbf{Title risk management:} a title change risks archive ambiguity
  (v65 carriers circulate under the old title). Resolution: the Data and
  Code Availability supersession sentence extended to retire the v65
  title explicitly, with ``no data, score, retention verdict, or
  reported number'' altered.
\item
  \textbf{Scope discipline:} no new numerals anywhere in the added prose
  (machine-checked: only intentional tokens ``4.6'' and ``65'' appear
  that were absent in v65); the six ``transferable elements'' are
  phrased as protocol design, not findings.
\item
  \textbf{Highlights kept within Elsevier convention} (five short
  bullets), mixing protocol and result bullets rather than grok's longer
  formulations.
\end{enumerate}

\subsection{4. Executed edits in v66
(anchors)}\label{executed-edits-in-v66-anchors}

\begin{enumerate}
\def\labelenumi{\arabic{enumi}.}
\tightlist
\item
  \textbf{Title} → ``Does Structural Elaboration Improve Biomass
  Forecasts? A Pre-Registered, Power-Aware Out-of-Sample Evaluation of
  Surplus-Production Models for Northern Cod (\emph{Gadus morhua})''
  (deepseek/qwen/grok option A).
\item
  \textbf{Highlights} --- five bullets: protocol, target alignment,
  empty retention, differential evidential weight, transferability.
\item
  \textbf{Abstract} --- ``The empirical conclusion applies\ldots{}'' +
  added protocol sentence + final sentence extended ``on
  decision-relevant quantities''.
\item
  \textbf{§1} --- new paragraph after shortcoming \#2: the two
  shortcomings are general; the case is a template for a
  forecast-evaluation protocol (cross-ref §2.3, §4.6).
\item
  \textbf{§2.3} --- generality paragraph after the retention-rule
  specification (declared comparators, frozen horizons, information-set
  documentation, rolling-origin scoring, uncertainty quantification,
  power calibration).
\item
  \textbf{§3.7} --- closing bookend sentence: a retention rule is an
  instrument with operating characteristics; non-retention is only as
  informative as the instrument's power under the relevant regime.
\item
  \textbf{§3.8} --- closing sentence: re-scoring on the rawest
  monitoring series is a general robustness discipline for
  reconstruction-based verdicts.
\item
  \textbf{§4.6 (new)} --- ``Broader Implications: A Transferable
  Protocol for Forecast Validation'': bounded empirical conclusion; six
  transferable elements; power-calibrated interpretation of negative
  results; collapse-window class-level misspecification as general
  regime-shift implication; conservative stress test (realized catches
  to M2--M4, retrospective reconstruction) with the no-overclaim hedge.
\item
  \textbf{§5 Conclusions} --- final paragraph: value of added structure
  must be demonstrated; protocol transferable to other stocks and
  ecological forecasting.
\item
  \textbf{Data and Code Availability} --- supersession extended to the
  v65 title-carrier.
\end{enumerate}

\textbf{Completeness statement (what was and was not checked):} every
claimed insert anchor was matched against v65 before application; all
ten edits applied; no empirically load-bearing sentence outside the ten
anchors was modified; numeral-diff machine check passed; PDFs verified
to contain the new title and §4.6 (27 pp, archive + submission
carriers). The SI (V7) was not changed --- nothing in it depends on the
framing edits; the cover letter was re-issued as v4 with the new title
only. The pending originality-response letter (20260919 v2) remains
valid and was not touched: its evidence concerns own-deposit overlap
sources, which is title-independent; if it is sent after v66 circulates,
one sentence noting the retitled superseding version may optionally be
added on review.

\end{document}
